% 第 46 章　量子力学的算符与矩阵语言
\chapter{量子力学的算符与矩阵语言}

前面三章给了我们一台完整的预测机器：量子态描述系统（第~43~章），薛定谔方程驱动它随时间变化（第~44~章），测量按玻恩规则吐出概率并把系统推进新态（第~45~章）。但这台机器还有一层我们一直没拆开的包装：量子态到底“是”一个什么数学对象？波函数 $\psi(x)$ 写得再好，也只是同一本书的一种抄本。本章要讲的是这台机器真正的内部语言——向量与矩阵。这不是数学家后来贴上去的装饰品：量子力学的全部怪异之处——测量结果为什么会离散、为什么有些量不能同时确定、为什么提问的顺序会改变答案——翻译进这门语言之后，全都变成中学生也能看懂的算术事实。学完本章，你再看“算符”“本征态”“对易”这些词，应该像看“速度”“加速度”一样亲切。

\step{反常识开场}

请先做一个谁都做过的动作：煮汤。先放盐再放糖，和先放糖再放盐，汤的味道没有区别。上班先回邮件再回消息，和反过来，结果也差不多。我们的整个日常经验都在反复确认一条“公理”：两件事情先后做的顺序，通常无所谓。数学也站在我们这边：$3\times5=5\times3$，先乘谁后乘谁都一样。

现在拿起你手边这本书，平放在桌上，封面朝上，书脊朝左。做两个动作。动作甲：把书绕左右方向的水平轴向前翻转九十度，让书立起来，封面朝向你。动作乙：把书绕竖直轴逆时针转九十度。先试“先甲后乙”，记下书的最终朝向；把书放回初始姿势，再试“先乙后甲”。

两个结果不一样。第一次，书立在那里，书脊朝下；第二次，书躺倒在桌面上，封面朝向侧面。同样的两个动作，同样的起点，只因为顺序不同，终点就不同。这不是魔术，任何三维物体都这样——可我们活了这么多年，居然很少注意到它。

先别急着说“这只是几何，跟物理有什么关系”。真正的悬念在这里：二十世纪二十年代，海森堡正是在显微镜下审视原子的光谱数据时发现，如果想让这些数对得上实验，描述原子的物理量必须遵守一种“顺序不同结果就不同”的乘法。他当时并不认识矩阵，是自己重新发明了一遍矩阵乘法，还为此惴惴不安，觉得物理学闯进了未知领域。一本普通的书早就天天向我们演示这种乘法了，只是没人把这两件事连起来。本章结束时你会看到：书的旋转、偏振片的顺序、位置和动量为什么不能同时测准，全是同一枚硬币的不同面。

\step{先猜一猜}

老规矩，先把你的预测写下来。

第一题。日常生活中“顺序无所谓”的操作和“顺序有所谓”的操作，你能不能各举三个例子？先给两个打样：煮汤先放盐还是先放糖，无所谓；穿袜子再穿鞋，和先穿鞋再穿袜子，差别可就大了。然后回答一个更难的问题：这两类操作之间，有没有一条干净的界线？是什么性质决定了一类操作可以随便交换顺序？

第二题。回忆第~45~章的三片偏振片实验：竖直片、水平片、斜四十五度片。把“先竖直再斜放”和“先斜放再竖直”两种顺序摆一摆——对单颗光子来说，这两种顺序的结局一样吗？如果你认为不一样，请回答：改变的到底是“光子被问的问题”，还是“光子本身”？

第三题，也是最尖锐的一题。假设物理量真的遵守“顺序不同结果不同”的乘法，也就是说存在这样的量 $A$ 和 $B$：$A$ 先作用、$B$ 后作用，不等于 $B$ 先作用、$A$ 后作用。你猜这会对“同时精确知道 $A$ 和 $B$”这件事造成什么后果？是毫无影响，还是暗藏杀机？写下你的直觉，读完本章再回来对答案。

\step{动手实验}

\begin{experiment}[一本书的两次旋转：顺序就是结果]
\textbf{材料}：一本平装书（或一部手机），一张桌面。

\textbf{第一步}：把书平放，封面朝上，书脊朝左。这是“初始状态”。

\textbf{第二步}：做动作甲——以书的左右方向为轴，把书向前翻转九十度，让书立起来，封面正对着你。接着做动作乙——以竖直方向为轴，把书逆时针转九十度。记下终点姿势：书立着，书脊朝哪边？封面朝哪边？

\textbf{第三步}：把书精确放回初始状态（这一步很关键，实验的公平性全靠它）。这次先做动作乙，再做动作甲。记下终点姿势。

\textbf{第四步}：对比两个终点。它们应当明显不同：一个书脊朝下立着，一个侧躺。如果看不出区别，检查你是不是两次都偷偷从同一个朝向起转的。

\textbf{记录与思考}：写下两句话——“甲后接乙”和“乙后接甲”各把书送到了什么姿势。然后回答：差别是谁造成的？是书不一样了吗？是动作变了吗？都不是。唯一的变量是顺序。这说明“旋转”这种东西，天然就不能用普通的数来记录——数从来不关心顺序，旋转关心。
\end{experiment}

有空的话，再做一个升级版：把其中一个九十度换成一百八十度，或者三个动作连做，你会发现顺序敏感不仅没有消失，反而花样更多。你正在亲手摆弄的东西，数学上叫“三维旋转群”，它和量子力学里自旋的数学结构是近亲——第~47~章我们还会回来走这门亲戚。

二维世界里有同一个现象的简化版，画在纸上就能看：一个不对称的三角形，“先旋转九十度再左右镜像”与“先镜像再旋转”的终点完全不同。

\begin{center}
\begin{tikzpicture}[scale=1.0]
% ---------- 第一行：先旋转再镜像 ----------
\begin{scope}
\draw (0,0) rectangle (1,1);
\fill[gray!40] (0,0) -- (1,0) -- (0,0.6) -- cycle;
\node at (0.5,-0.35) {\footnotesize 初始};
\end{scope}
\begin{scope}[xshift=1.8cm]
\draw (0,0) rectangle (1,1);
\fill[gray!40] (1,0) -- (1,1) -- (0.4,0) -- cycle;
\node at (0.5,-0.35) {\footnotesize 旋转$90^\circ$};
\end{scope}
\begin{scope}[xshift=3.6cm]
\draw (0,0) rectangle (1,1);
\fill[gray!40] (0,0) -- (0,1) -- (0.6,0) -- cycle;
\node at (0.5,-0.35) {\footnotesize 再镜像};
\end{scope}
% ---------- 第二行：先镜像再旋转 ----------
\begin{scope}[xshift=5.6cm]
\draw (0,0) rectangle (1,1);
\fill[gray!40] (0,0) -- (1,0) -- (0,0.6) -- cycle;
\node at (0.5,-0.35) {\footnotesize 初始};
\end{scope}
\begin{scope}[xshift=7.4cm]
\draw (0,0) rectangle (1,1);
\fill[gray!40] (1,0) -- (0,0) -- (1,0.6) -- cycle;
\node at (0.5,-0.35) {\footnotesize 镜像};
\end{scope}
\begin{scope}[xshift=9.2cm]
\draw (0,0) rectangle (1,1);
\fill[gray!40] (1,1) -- (1,0) -- (0.4,1) -- cycle;
\node at (0.5,-0.35) {\footnotesize 再旋转$90^\circ$};
\end{scope}
% 箭头
\draw[-{Stealth}] (1.15,0.5) -- (1.65,0.5);
\draw[-{Stealth}] (2.95,0.5) -- (3.45,0.5);
\draw[-{Stealth}] (6.75,0.5) -- (7.25,0.5);
\draw[-{Stealth}] (8.55,0.5) -- (9.05,0.5);
\end{tikzpicture}
\end{center}

\step{旧模型遇到麻烦}

整个经典物理学都建立在一个从来不必明说的记账习惯上：每个物理量就是一个数。位置是三个数，速度是三个数，能量是一个数。数的好处是运算干净：加法交换，乘法交换，$a+b=b+a$，$ab=ba$。测量也只是“把那个数读出来”，读数的顺序当然无所谓——你先称体重再量身高，和先量身高再称体重，得到的是同一份体检表。

麻烦从两个方向挤过来。

第一个方向就是刚才的实验。旋转是实实在在的物理操作，可是“先做旋转甲再做旋转乙”这件事，用普通的数根本记不了账——数没有顺序概念。十九世纪的数学家为此发明了矩阵：一张数字的表格，作用在向量上；两张表格相乘代表“接连作用”，而这种乘法天然不交换。矩阵不是为量子力学发明的，它是被几何逼出来的。

第二个方向才是物理学的地震。第~45~章里我们已经亲眼看到：对光子先问“竖直方向偏振吗”再问“四十五度方向偏振吗”，与颠倒提问顺序，是两次不同的实验，结局分布都不同。测量顺序进入了物理本身。经典记账法在这里彻底破产：如果每个物理量只是一个预先写好的数，测量只是抄数，那么先抄哪一笔根本不该有区别。唯一的出路是承认：量子系统的“物理量”不是数，而是某种对状态施加的操作——而操作，就像书的旋转一样，可以有顺序敏感性。

还有第三处裂缝，为后来的不确定性埋下伏笔。在经典力学里，位置和动量各自是一个干干净净的数，想知道多少精度都可以，互不干扰。可实验一再表明：把电子的位置限制得越精确，它的动量分布就散得越开。第~43~章的波包已经演示过这一点。如果位置和动量只是两列普通的数，这种此消彼长毫无道理。但如果它们是两个“不可交换的操作”，这一切将自动成为算术推论——这正是本章“数学翻译器”一节要做的事。

\step{发明新概念}

现在像海森堡、薛定谔、狄拉克当年那样，把新账簿一页一页设计出来。我们要发明的概念只有四个，每一个都先给直觉，再给名字。

\textbf{第一页：状态是向量，但不一定指向空间中的方向。}你早就用过这种向量，只是没叫它向量。一张期末成绩单：语文八十五，数学九十二，英语七十八——三个数排成一列，这就是一个“成绩向量”。它不指向东也不指向北，它指向“成绩空间”里的一个位置。再举两个：调色软件里一种颜色由红、绿、蓝三个分量确定，是一个“颜色向量”；你的月度账单上餐饮、交通、房租各项开支排成一列，是一个“开销向量”；音乐软件里一首歌也可以是一列数——节奏、响度、人声比例、低音强度各占一格，推荐算法“猜你喜欢”干的正是向量之间的加减与比较，你每天在消费的歌单背后就是本章的数学。向量的本质不是“箭头”，而是“一组合格的坐标，遵守合理的加减规则”。这个比喻要立刻标明边界：成绩单向量像量子态的地方是，都用一列数完整描述某种状态；不像的地方有两条——量子态的分量一般是复数，而成绩是实数；更关键的是，测量会改变量子态，而查看成绩单不会改变成绩。量子力学的第一条大举动，就是把“系统此刻是什么状态”的答案写成一个抽象向量，称为\textbf{态矢量}。波函数 $\psi(x)$ 只是态矢量在“位置基底”下的坐标写法，就像同一个地点既能写成经纬度、也能写成“钟楼往北三百米”。

\textbf{第二页：基底是你选定的报数方式。}同一个地点，用经纬度报是一对数，用“距钟楼多远、偏东多少度”报是另一对数。地点没变，坐标变了，因为报数的“参照框架”——也就是\textbf{基底}——换了。量子态也一样：同一个光子偏振态，可以用“竖直、水平”这对基来报坐标，也可以用“四十五度、一百三十五度”那对基来报。坐标数字不同，态是同一个。这件事在量子力学里威力巨大：所谓“选择测量什么”，数学上就是“选择用哪套基底来展开态矢量”。第~45~章说的“你选定一个问题”，现在有了精确的对应物。

\begin{center}
\begin{tikzpicture}[scale=1.1,>=Stealth]
% 标准基底
\draw[->] (0,0) -- (2.6,0) node[right] {$e_1$};
\draw[->] (0,0) -- (0,2.6) node[above] {$e_2$};
% 旋转基底
\draw[->,thick,blue!60!black] (0,0) -- (30:2.8) node[right] {$e'_1$};
\draw[->,thick,blue!60!black] (0,0) -- (120:2.8) node[above] {$e'_2$};
% 向量 v
\draw[->,very thick,red!70!black] (0,0) -- (0.6,1.6) node[above right] {$v$};
% 在标准基上的投影
\draw[dashed,gray] (0.6,1.6) -- (0.6,0);
\draw[dashed,gray] (0.6,1.6) -- (0,1.6);
% 在旋转基 e'_1 上的投影（垂足约 (1.14,0.66)）
\draw[dashed,blue!60!black] (0.6,1.6) -- (1.14,0.66);
\node at (3.1,2.4) {\footnotesize 同一支箭头，两套坐标};
\end{tikzpicture}
\end{center}

\textbf{第三页：算符是对状态施加的操作。}修图软件是最好的直觉入口。一张照片是一个状态；“旋转九十度”“左右镜像”“调亮”“磨皮”，每一个都是把旧状态变成新状态的操作——这就是\textbf{算符}。注意两件小事。第一，算符关心顺序：先旋转九十度再左右镜像，和先镜像再旋转，照片结果不一样——你不妨在手机上试试，这和书的旋转是同一个现象。第二，量子力学要求算符是\textbf{线性}的：作用在两个态的叠加之和上，等于分别作用再相加。这个要求对应一个深刻的物理事实——第~43~章的叠加原理：量子态可以像波一样叠加，干涉项一个都不能丢。调亮是线性操作的好榜样（两张照片叠加后再调亮，等于各自调亮再叠加），而“裁剪成正方形”就不是（叠加后的照片和各自裁剪的结果对不上）。这个比喻也要划线：滤镜像算符的地方是“对输入做操作、关心顺序”；不像的地方是，量子算符作用完状态本身仍是确定的，随机性只出现在测量那一步——别把算符和测量混为一谈。

\textbf{第四页：本征态是操作的“不动方向”。}绝大多数操作会改变状态。但每个操作都偏爱一些特殊的状态：作用上去之后，状态的方向纹丝不动，只是整体被拉长、缩短、翻个面，顶多再乘上一个相位。比如“绕竖直轴旋转”这个操作，对一支竖直放置的箭毫无作为——箭就是它自己的本征态，本征值是 $1$（完全不变）。“左右镜像”对一张本来就左右对称的照片也只是保持原样。这种“操作的特殊状态”叫\textbf{本征态}，被乘上的那个倍数叫\textbf{本征值}。为什么量子力学对它念念不忘？因为第~45~章的账簿里写着：测量某个量时，可能出现的确定结果，恰恰是这个量对应算符的本征值；测量后系统落入的新态，恰恰是对应的本征态。原子光谱为什么是一条条离散的亮线，而不是连续彩虹？因为原子能量算符的本征值清单是离散的。离散性不再是谜，它是“操作的不动方向只有有限几个倍数”这条算术事实的物理化身。

把四页账簿合起来，量子力学的核心句就成形了：态是向量，测量量是算符，测量可能结果是本征值，测后状态是本征态。两个算符能不能交换顺序——它们的\textbf{对易}性质——决定这两个量能不能同时拥有确定值。

\step{一句话模型}

\begin{oneline}
量子力学把“系统的状态”记成一个抽象向量，把“可观测量”记成作用在向量上的操作；每个操作都有自己偏爱的“不动方向”（本征态）和对应的倍数（本征值），测量只能报出这些倍数之一；两个操作若不能交换顺序，对应的两个量就不能同时拥有确定值——量子世界的全部“怪异”，都是这门向量算术的直译。
\end{oneline}

这个模型像什么？像一本用坐标纸记账的流水账：换一套坐标轴，数字全变，账本身没变。它不像什么？不像一叠写着标准答案的卡片：账上记的不是“系统藏着哪些预先定好的数”，而是“每种操作会把状态推向哪里”。也请再次标明它不是什么：态矢量不是空间里一支真实的箭头，算符不是一台真实机器，它们是对“制备—操作—测量”这一整套实验规程的数学压缩。

\step{数学翻译器}

下面把直觉逐条翻译成符号。主线层只需要你会做“表格乘箭头”的算术；更远的推导放在数学桥和挑战层。

\textbf{翻译一：向量与坐标。}选定基底后，一个态写成一列数，例如二分量情形
\[
|\psi\rangle=\begin{pmatrix} a \\ b \end{pmatrix},
\]
其中 $a$、$b$ 可以是复数，分别叫“沿第一个基矢、第二个基矢的振幅”。第~43~章的规则照旧：测量对应基矢所指的那个结果，概率是振幅模的平方，即 $|a|^2$ 与 $|b|^2$，所以合格的态必须满足 $|a|^2+|b|^2=1$——概率加起来是一。立刻检查特殊情况：若 $b=0$，态就是第一个基矢本身，测量结果百分之百是第一个——坐标 $(1,0)$ 与“本征态”恰好是同一回事，自洽。

\textbf{翻译二：算符是数字表格，乘法是“接连作用”。}同样二分量的情形，一个算符写成 $2\times2$ 的表格
\[
\hat{A}=\begin{pmatrix} p & q \\ r & s \end{pmatrix},
\]
它作用在向量上就是表格乘箭头，规则是“逐行点乘”：结果的上一格是 $pa+qb$，下一格是 $ra+sb$。两个算符接连作用，对应两张表格相乘；而这种乘法不保证交换，$\hat{A}\hat{B}$ 与 $\hat{B}\hat{A}$ 可以是两张不同的表格。这在数学桥上我们算一个真例子。

\textbf{翻译三：本征方程。}“操作不改变方向，只放大倍数”写成一行：
\[
\hat{A}\,|\psi\rangle=\lambda\,|\psi\rangle,
\]
读作：算符 $\hat{A}$ 作用在它自己的本征态 $|\psi\rangle$ 上，等于原态乘上本征值 $\lambda$。用零、一、反转三种特殊情况检查它。若 $\lambda=1$：操作对该态完全透明，如旋转轴上的箭。若 $\lambda=0$：操作把这个态“消灭”成零，例如竖直偏振片遇到水平偏振光——全挡，这正好解释第~45~章十字交叉的两片为何全黑。若 $\lambda=-1$：方向反转，对应镜像对称中“反对称”的态，如把左右互换后整体变号的驻波。三种情形都在实验里找得到影子。

\textbf{翻译四：对易子与不确定性。}两个算符顺序之差的度量叫\textbf{对易子}：
\[
[\hat{A},\hat{B}]=\hat{A}\hat{B}-\hat{B}\hat{A}.
\]
若它等于零，两个操作可交换，存在一套共同的本征态，两个量能同时有确定值——先测谁后测谁都一样。若它不等于零，麻烦来了。量子力学的地基里埋着一条从实验与理论双向焊死的关系：位置算符 $\hat{x}$ 与动量算符 $\hat{p}$ 满足
\[
[\hat{x},\hat{p}]=i\hbar,
\]
其中 $\hbar\approx1.05\times10^{-34}\,\mathrm{J\cdot s}$ 是约化普朗克常数，$i$ 是虚数单位。右边不是零——位置和动量永远不可交换。从这一条可以严格推出海森堡不确定关系 $\Delta x\,\Delta p\geq\hbar/2$（推导见挑战层）。请特别注意这条结论的性质：它没有提到任何仪器。不确定关系不是“测量手段太粗糙”的道歉，而是“$\hat{x}$ 与 $\hat{p}$ 不交换”这条结构事实的影子——哪怕有一台完美仪器，结果也不会变。

\textbf{翻译五：算一笔真实数量的账。}不确定关系在日常世界为什么不露面？代入数字看。一粒质量约 $10^{-4}\,\mathrm{kg}$ 的细尘（花粉颗粒的量级），假设我们把它的位置确定到头发丝量级，$\Delta x\sim10^{-4}\,\mathrm{m}$，那么动量不确定度至少约 $\hbar/(2\Delta x)\sim5\times10^{-31}\,\mathrm{kg\cdot m/s}$，折合速度不确定度约 $5\times10^{-27}\,\mathrm{m/s}$——以这个速度，尘埃漂移一根头发丝的宽度需要比宇宙年龄还长的时间，完全没有观测后果。换一个极端：把电子限制在原子尺度内，$\Delta x\sim10^{-10}\,\mathrm{m}$，电子质量约 $9.1\times10^{-31}\,\mathrm{kg}$，速度不确定度立刻冲到约 $6\times10^{5}\,\mathrm{m/s}$，与原子中电子的特征速度同量级——不确定性直接主宰原子的尺寸与稳定性。同一个公式，两个世界，分野全在 $\hbar$ 那小小的 $10^{-34}$ 上。再检查一个极限：若形式上令 $\hbar\to0$，对易子归零，一切算符恢复交换，量子力学自动退回经典记账——公式的边界行为正确。

\textbf{翻译六：平均也是一笔算得死的账。}单次测量是随机的，但“完全相同的制备重复一万次，结果的平均值是多少”却是理论能精确预言的，这个平均值叫\textbf{期望值}：把本征值清单上的每个数按它的概率加权求和。这和掷骰子平均得 $3.5$ 是同一套算术，只是骰子的清单是 $1$ 到 $6$，量子系统的清单是算符的本征值。立刻做特殊检查：若系统的态恰好是某本征态，对应结果的概率是 $1$、其余全是 $0$，加权和就退化为那一个本征值——确定的态，期望值就是它自己，涨落为零；这正是“本征态对应确定测量结果”的另一种说法，账目自洽。期望值不是单次测量会报出的数——骰子永远掷不出 $3.5$——它是大量重复的集体性质，这一点和第~45~章“理论只承诺比例”的立场完全一致。

\begin{mathbridge}
\textbf{一个真的能算完的例子。}取两张 $2\times2$ 表格：
\[
\hat{X}=\begin{pmatrix} 0 & 1 \\ 1 & 0 \end{pmatrix},\qquad
\hat{Z}=\begin{pmatrix} 1 & 0 \\ 0 & -1 \end{pmatrix}.
\]
按“逐行点乘”的规则亲手算两笔（值得花五分钟）：
\[
\hat{X}\hat{Z}=\begin{pmatrix} 0 & -1 \\ 1 & 0 \end{pmatrix},\qquad
\hat{Z}\hat{X}=\begin{pmatrix} 0 & 1 \\ -1 & 0 \end{pmatrix},
\]
两者恰好差一个负号：$\hat{X}\hat{Z}=-\hat{Z}\hat{X}$，对易子 $[\hat{X},\hat{Z}]=2\hat{X}\hat{Z}\neq0$。顺序在这里不仅“有关系”，而且是反着号的关系。再求 $\hat{Z}$ 的本征态：设 $\hat{Z}\begin{pmatrix}a\\b\end{pmatrix}=\lambda\begin{pmatrix}a\\b\end{pmatrix}$，逐行写出 $a=\lambda a$、$-b=\lambda b$，只有两条出路——$b=0,\lambda=1$ 或 $a=0,\lambda=-1$。也就是说 $\hat{Z}$ 的本征态正是两个基矢 $(1,0)$ 与 $(0,1)$，本征值 $+1$ 与 $-1$。同理可求 $\hat{X}$ 的本征态是 $(1,1)/\sqrt2$ 与 $(1,-1)/\sqrt2$——用 $\hat{Z}$ 的基底报数，它们是两个基矢的等权叠加。于是“$\hat{X}$ 有确定值的态，对 $\hat{Z}$ 而言是两个结果各半概率”——不可交换与“不能同时确定”在纸上直接见了面。这两张表格不是玩具：它们就是著名的泡利矩阵中的两张，第~47~章讲自旋时它们会作为主角登场。
\end{mathbridge}

\begin{challenge}
\textbf{从对易关系到不确定关系，再到无穷维的暗礁。}对任意两个算符与任意态，定义均方偏差 $\Delta A$、$\Delta B$，用施瓦茨不等式作用于向量 $(\hat{A}-\langle\hat{A}\rangle)|\psi\rangle$ 与 $(\hat{B}-\langle\hat{B}\rangle)|\psi\rangle$，可推出普遍不等式
\[
\Delta A\,\Delta B\geq\tfrac12\,\bigl|\langle[\hat{A},\hat{B}]\rangle\bigr|.
\]
代入 $[\hat{x},\hat{p}]=i\hbar$，右边给出 $\hbar/2$，海森堡关系到手——它是一条纯粹的向量几何定理，与仪器无关。另一个暗礁：位置、动量这类算符生活在无穷维空间里，本征值构成连续谱，“本征态”不再是普通向量而是分布（狄拉克 $\delta$），于是 $\hat{x}$ 与 $\hat{p}$ 的定义域不能覆盖整个空间，$[\hat{x},\hat{p}]=i\hbar$ 严格说只在公共定义域上成立。本章所有结论在物理上不受影响，但“每个对称算符都有好本征态”这类天真断言在无穷维需要修补——这是泛函分析的地界，留给你将来的某本书。
\end{challenge}

\step{模型的边界}

每一套漂亮的语言都有自己的护城河与沼泽，算符语言也不例外。

\textbf{边界一：哈密顿算符不自动等于“总能量”。}第~44~章里，薛定谔方程 $i\hbar\,\partial_t|\psi\rangle=\hat{H}|\psi\rangle$ 用哈密顿算符 $\hat{H}$ 生成时间演化。在许多常见情形——封闭系统、约束力不做功、势能不显含时间——$\hat{H}$ 的本征值正好就是系统的总能量，于是“能量算符”这个叫法深入人心。但这不是定义，是有条件的结论。系统在随时间变化的外部装置中（比如被一台按程序调制的磁场驱动），$\hat{H}$ 可以显含时间，本征值不再对应一个守恒的“总能量”；在磁场中的带电粒子那里，出现在 $\hat{H}$ 里的是正则动量而不是机械动量，“动能加势能”的朴素读法也会失灵。准确的说法是：量子力学是以哈密顿算符生成演化的理论。有些流行读物讲的“哈密顿量子力学”并不是独立于薛定谔、海森堡表述之外的第三套理论，它只是强调“演化由 $\hat{H}$ 生成”这一共同的骨架。

\textbf{边界二：薛定谔绘景与海森堡绘景是同一场电影的两种放映机。}薛定谔绘景让态矢量随时间转、算符静止；海森堡绘景反过来，把全部时间演化塞进算符，态一动不动。两者的全部可观测预言逐条相同——期望值、概率、能谱，没有一项实验能区分它们。这就像描述钟摆：你可以盯着一个摆动的小球记录它的位置随时间变化（态在动），也可以让小球保持“初始标签”不动、让“此刻位置”这把尺子随时间变形（算符在动）。电影是同一部。哪种绘景顺手就用哪种：算能谱常用前者，找守恒量、做微扰常用后者——后者里藏着一个一眼可见的漂亮结论：一个算符若与哈密顿算符对易，它就不随时间变化，于是“守恒量”翻译成新语言就是“与 $\hat{H}$ 对易的量”，能量、动量、角动量守恒在同一条规则下站成了一排。把两种绘景当成“两种量子力学”是初学时常见的误会。

\textbf{边界三：矩阵力学与波动力学是同一理论的两种坐标。}海森堡的矩阵、薛定谔的波函数，看似两门手艺，薛定谔本人于1926年证明了它们数学等价：波函数 $\psi(x)$ 无非是态矢量在位置基底下的坐标，微分算符 $-i\hbar\,\partial_x$ 无非是动量算符在同一基底下的长相。选基底如选地图投影，改变不了地球。

\textbf{边界四：算符语言不替你选诠释。}本征值清单、对易关系、演化方程，都是各家诠释（第~45~章末尾的那张中立地图）共同承认的硬核。但“测量瞬间态如何更新”在方程里仍是一条外加规则，退相干能解释经典表象如何浮现，却是否单独解决了测量问题仍有争议。学会矩阵语言，你拿到的是各家争论的共同赛场，而不是裁判的哨子。

\textbf{边界五：别把复数当玄学。}态分量是复数，可一切可观测的预言——概率、期望值、本征值——都是实数（这正是“可观测量对应厄米算符”的算术内容）。复数在这里的角色更像交流电计算里的相量：中间步骤用相位记账，最后报数时取实的结果。

\step{回头解释开场现象}

现在回到桌上那本书。两次旋转顺序不同、终点不同，这个看似琐碎的现象，在本章语言里获得了一个精确的身份：三维旋转操作对应一组不可交换的算符（数学上它们是转动群的生成元），“先甲后乙”与“先乙后甲”对应两张不同的矩阵乘积。你的书从始至终在演示矩阵乘法，只是没人给你发那张乘法表。

再往下挖一层。量子力学并没有“发明”顺序敏感，它只是在微观世界里撞见了同一结构，并且撞得更深。对普通的数，$AB-BA=0$；对书的旋转，$AB-BA$ 是一个小角度旋转；对位置与动量，$[\hat{x},\hat{p}]=i\hbar$，一个虽然微小却永远不为零的常数。三条语句排在一起，一条清晰的世界观浮现出来：自然界的基本对象不是“带着数值的标签”，而是“彼此作用的操作”；数只是操作在特殊情况下（可交换时）退化成的影子。经典物理学之所以几百年没发现这一点，纯粹因为 $\hbar$ 太小——第“数学翻译器”一节那笔尘埃的账就是证明：在宏观尺度上，不可交换性被压缩到小数点后二十多位，粗看一切交换，细看全是矩阵。

至于第二道猜题——偏振片顺序——现在也是算术：偏振沿不同方向的投影算符彼此不可交换，交换两个投影等于交换两次测量，也就等于换了一个实验。第三道猜题的答案你已经亲手在数学桥里算出来了：不可交换的算符没有共同的全部本征态，一个量的确定态在另一个量看来必然是若干结果的叠加——“不能同时精确知道”不是实验技术宣言，是 $2\times2$ 表格的乘法性质。

这套语言的力量远不止解释旧现象。它已经是工程图纸：量子计算机的每一个量子比特就是一个二分量态矢量，每一个逻辑门就是一张酉矩阵——$X$ 门实现“翻转”，$H$ 门把确定态变成等权叠加，整块量子芯片就是一本巨大的矩阵乘法练习册；核磁共振与医学磁共振成像的整套脉冲序列设计，全部用自旋算符与它们的指数演化写成；原子和分子的能级计算，本质就是求一个巨大矩阵的本征值清单。向量与矩阵从描述语言变成了生产工具。

\step{脑洞实验与挑战题}

\begin{thought}[如果 $\hbar$ 大三十五亿亿亿倍]
对易关系 $[\hat{x},\hat{p}]=i\hbar$ 里，$\hbar$ 的日常渺小是经典世界存在的全部理由。做一个极端尺度的思想实验：设想某个平行宇宙里 $\hbar$ 不是 $10^{-34}\,\mathrm{J\cdot s}$，而是约 $1\,\mathrm{J\cdot s}$。那么对一只 $0.1\,\mathrm{kg}$ 的苹果，把位置确定到一厘米，速度不确定度至少是几米每秒——你“看清”苹果放在哪儿的那一刻，它就糊成了一片速度的云雾；想拍清楚苹果飞得多快，它又在桌上摊开成一摊位置的迷雾。再往下想：这台宇宙里还能有“轨道”吗？行星绕着恒星，既要有确定的位置（在轨道某处）又要有确定的动量（沿轨道切向）——两个要求在那个宇宙里互相封杀，开普勒椭圆根本画不出来，取而代之的是一团弥散的“行星云”。我们的宇宙之所以有桌子、轨道和天气预报，不是因为我们幸运，而是因为 $\hbar$ 恰好小到让 $[\hat{x},\hat{p}]$ 在宏观尺度上约等于零。常数的大小，决定了什么种类的世界能够存在。
\end{thought}

\begin{puzzles}
\item \textbf{共同本征态的存在性。}算符 $\hat{A}$ 的本征态是“沿 $x$ 轴的箭”，算符 $\hat{B}$ 的本征态是“沿四十五度方向的箭”。假设 $\hat{A}$ 与 $\hat{B}$ 可以交换。试论证：它们至少存在一套共同的本征态基底。反过来，若两者的本征方向处处对不上，你能对 $[\hat{A},\hat{B}]$ 下什么结论？（思路提示：取 $\hat{A}$ 的一个本征态 $|a\rangle$，考察 $\hat{A}(\hat{B}|a\rangle)$，利用交换性把 $\hat{A}$ 换到最前面；再想“对不上”意味着什么。）
\item \textbf{顺序的代价。}一个光子先通过竖直偏振片，再遇到“沿四十五度方向提问”的装置；另一个一模一样的光子，顺序对调。已知竖直偏振态用四十五度基底展开是等权叠加。试比较两个光子最终“通过四十五度检验”的概率，并说明：交换顺序之所以改变了结局，究竟改了哪一步？（思路提示：概率只在“用新基底报数”那一步出现；注意第一次测量后态已经更新，第~45~章的账簿要翻出来用。）
\item \textbf{两个“纯读取”的操作为什么不交换。}直觉说：测量身高和测量体重互不影响，所以“先测 $A$ 再测 $B$”与顺序无关应该普遍成立。请用本章语言指出这个推理暗含了什么假设，并举一个反例说明该假设在量子层面的具体死法。（思路提示：暗含的假设是“所有算符共享同一套本征态且测量不改态”；回想数学桥里 $\hat{X}$ 的本征态在 $\hat{Z}$ 基底下的长相。）
\item \textbf{两种放映机。}一个态在薛定谔绘景里随时间旋转，在海森堡绘景里静止。若有人声称“海森堡绘景里系统不演化，所以是错的”，请反驳他，并说明哪种实验数据可以——以及不可以——仲裁两种绘景之争。（思路提示：可观测的只有期望值与概率；检查两者在两种绘景下的表达式，注意“谁动”只是记账方式。）
\end{puzzles}

本章把量子力学从“波的故事”升级成了“向量与操作的故事”。下一站是一个躲不掉的考验：自旋。它是最纯粹的二分量系统，泡利矩阵是它的母语，而你刚刚学会的所有算术——基底、本征值、不可交换——都将在那里得到一次不容讨价还价的实战检验。
